\clearpage
\SourceLeaf{279}
\renewcommand{\BookRunningTitle}{De nexu inter multitudinem classium}
\renewcommand{\Article}[1]{\par\medskip\Needspace{4\baselineskip}\pdfbookmark[1]{Art. #1}{nexus-prior-art-#1}\centerline{\textbf{#1.}}\nobreak\smallskip}
\pdfbookmark[0]{De nexu inter multitudinem classium. I}{nexus-prior}
\begin{center}
[I.]\par\medskip
\Large\bfseries DE NEXU INTER MULTITUDINEM CLASSIUM, IN QUAS FORMAE BINARIAE SECUNDI GRADUS DISTRIBUUNTUR, EARUMQUE DETERMINANTEM.\par\medskip
\normalsize COMMENTATIO PRIOR\par\medskip
SOCIETATI REGIAE EXHIBITA 1834\ .\ .\ .\ .
\end{center}
\BookRule{.18\textwidth}

\Article{1}
Triginta tres iam elapsi sunt anni, ex quo principia nexus mirabilis, cui haec commentatio dicata est, deteximus, uti iam in fine \emph{Disquisitionum Arithmeticarum} annunciatum est. Sed aliae occupationes ab hac scrutatione per longum tempus detraxerant, donec recentiori tempore ad eam reverti et per novas curas eam ampliare contigit. Attamen quum haec nova Arithmeticae Sublimioris pars limites unius commentationis excedat, haecce prior formis determinantium negativorum dicata erit: formae vero determinantium positivorum, quae tractationem prorsus peculiarem requirunt, commentationi alteri reservatae manere debebunt.

\Article{2}
Basis totius argumenti est disquisitio peculiaris circa multitudinem omnium combinationum valorum integrorum, quos duo numeri integri indefiniti $x,y$ intra ambitum praescriptum accipiunt. Manifesto hoc problema etiam sub aspectu geometrico exhiberi potest, ut eruatur multitudo \emph{numerorum complexorum}, quorum repraesentatio intra figuram praescriptam cadit. Indoles figurae ex indole lineae quae eam circumdat, adeoque pendebit vel ab unica aequatione inter coordinatas $x,y$ (quoties peripheria est curva in se rediens) vel a pluribus huiusmodi 
\SourceLeaf{280}
aequationibus (quoties .constat e pluribus partibus curvis seu rectis), pendebitque ab arbitrio nostro, utrum puncta numeris integris complexis respondentia, si quae forte in ipsa peripheria sint, multitudini annumerare velimus an inde excludere.

In repraesentatione analytica problematis conditiones illius limitationis semper ita exhiberi poterunt, ut functio data variabilium $x,y$ vel una vel plures $P,Q,R$ etc. nancisci debeant valores positivos, vel non-negativos (prout valor $0$ vel excluditur vel admittitur).

Ita e. g. si figura praescripta est circulus, cuius radius $=\surd A$, dum centrum cadit in punctum numero complexo integro respondens, conditio analytica erit, ut $A-xx-yy$ non sit negativus, siquidem, quod semper supponemus, puncta in ipsa peripheria sita retinere placet. Si figura est triangulum, tres functiones lineares $ax+by+c$, $a'x+b'y+c'$, $a''x+b''y+c''$ valores non-negativos habere debent, similiterque in aliis casibus.

\Article{3}
Solutio problematis \emph{exacta}, generaliter loquendo, ita procedere debet, ut primo e natura conditionum variabilis altera e. g. $x$ intra limites coërceatur, inter quos valores singuli integri deinceps percurrant, et quot valores integri alterius $y$ singulis respondeant, eruere oportet, quorum multitudines dein in summam colligi debent. In casibus specialibus plerumque aderunt artificia specialia ad laborem abbreviandum.

E. g. si figura, ut supra, est circulus, cuius radius $=\surd A$, sit $r$ integer proxime minor quam $\surd A$, vel ipse $\surd A$, si $A$ est quadratum. Perinde sint $r',r'',r'''$ etc. $r^{(r)}$ integri proxime minores quam $\surd(A-1)$, $\surd(A-4)$, $\surd(A-9)$ etc. usque ad $\surd(A-rr)$. Tunc multitudo quaesita erit
\[
\begin{aligned}
&=2r+1+2(2r'+1)+2(2r''+1)+2(2r'''+1)+\text{etc.}\\
&=1+4r+4r'+4r''+4r'''+\text{etc.}+4r^{(r)}
\end{aligned}
\]

Brevior erit in hoc exemplo methodus sequens. Sit $q$ integer proxime minor quam $\surd\tfrac12 A$ (vel huic aequalis, quoties est integer), atque $r^{(q+1)}$, $r^{(q+2)}$, $r^{(q+3)}$ etc. integri proxime minores quam $\surd(A-(q+1)^2)$, $\surd(A-(q+2)^2)$, $\surd(A-(q+3)^2)$ etc. usque ad $\surd(A-rr)$. Tunc erit multitudo quaesita
\[
=4qq+1+4r+8(r^{(q+1)}+r^{(q+2)}+r^{(q+3)}+\text{etc.}+r^{(r)})
\]

\SourceLeaf{281}
Per hanc formulam eruta est multitudo
\[
\begin{array}{r|r||r|r||r|r}
A&&A&&A&\\\hline
100&317&1000&3149&10000&31417\\
200&633&2000&6293&20000&62845\\
300&949&3000&9425&30000&94237\\
400&1257&4000&12581&40000&125629\\
500&1581&5000&15705&50000&157093\\
600&1885&6000&18853&60000&188453\\
700&2209&7000&21993&70000&219901\\
800&2521&8000&25137&80000&251305\\
900&2821&9000&28269&90000&282697\\
1000&3149&10000&31417&100000&314197
\end{array}
\]

\Article{4}
Ad propositum nostrum non requiritur determinatio exacta, sed potius indagatio expressionis, quae ad multitudinem exactam quam prope velis accedere potest, dum limites in infinitum ampliantur. Sed ante omnia quum haec aliquid vagi involvant, rem exactius explicare oportet.

Supponemus itaque, functiones $P,Q,R$ etc. praeter variabiles $x,y$ implicare elementum constans $k$, ita ut singulae $P,Q,R$ etc. sint functiones homogeneae trium quantitatum $x,y,k$. Hoc pacto figura per aequationes $P=0,Q=0,R=0$ etc. determinata pendebit a $k$, ita ut valoribus diversis ipsius $k$ respondeant figurae similes et respectu initii coordinatarum similiter positae, dimensionesque lineares similes valoribus ipsius $k$, areae valoribus ipsius $kk$ proportionales erunt. Denotetur iam multitudo punctorum intra figuram per $M$, area per $V$, patetque $M$ et $V$, crescente $k$, crescere debere; crescente vero $k$ in infinitum, $M$ et $V$ ad rationem aequalitatis quam proxime velis accedent, vel si elementarem claritatem postulas, proposita quantitate quantumvis parva $\lambda$, semper assignari poterit terminus talis, ut pro quolibet valore ipsius $k$ hunc terminum superante certo $\frac MV$ iacere debeat inter $1-\lambda$ et $1+\lambda$. Secundum morem suetum hoc ita indicare licet: fieri $M=V$ pro valore infinito ipsius $k$.

In exemplo nostro conditio requisita locum tenet, statuendo $k=\surd A$, curvaque fit circulus, cuius area $=\pi A$, denotante $\pi$ semicircumferentiam circuli pro radio $=1$. Numeri supra traditi convergentiam luculenter addigitant.

\SourceLeaf{282}
Ceterum si operae pretium esset, facile demonstrationem illius theorematis antiquo rigore absolvere possemus, quam tamen hocce quidem loco supprimere maluimus ad difficiliora properantes.

\Article{5}
In hacce commentatione limes per \emph{unicam} aequationem talem exprimetur $axx+2bxy+cyy=A$, ita quidem ut $a,b,c$ sint integri, atque $bb-ac$ numerus negativus quem statuemus $=-D$. Manifesto curva figuram definiens erit ellipsis, patetque facile, quadrata semiaxium esse radices aequationis
\[
(ac-bb)qq-(a+c)Aq+AA=0\quad\text{sive}\quad
=A\left(\frac{a+c\pm\surd(4bb+(a-c)^2)}{2(ac-bb)}\right)
\]

Productum harum radicum fit $\frac{AA}{ac-bb}=\frac{AA}D$, proin area ellipsis $=\frac{\pi A}{\surd D}$. Hinc itaque colligitur, multitudinem omnium combinationum valorum integrorum ipsarum $x,y$, pro quibus $axx+2bxy+cyy$ valorem $A$ non superet, crescente $A$ continuo magis appropinquare ad $\frac{\pi A}{\surd D}$, et pro $A$ infinito huic valorem aequalem statui debere. Ceterum manifestum est, hocce respectu nihil interesse, utrum combinatio $x=0,y=0$ reliquis annumeretur, an inde excludatur. Hoc itaque modo multitudo quaesita (in ratione posteriori) nihil aliud est, nisi aggregatum multitudinum repraesentationum singulorum numerorum $1,2,3\ldots A$ per formam binariam secundi gradus $axx+2bxy+cyy$; et quum inter illos numeros alii omnino per hanc formam repraesentari nequeant, alii plures, alii pauciores repraesentationes admittant, quantitas $\frac\pi{\surd D}$ consideranda erit tamquam valor medius multitudinis repraesentationum numeri positivi indefiniti per formam quamlibet, cuius determinans $=-D$.

